% !TEX program = pdflatex
\documentclass[12pt, a4paper, bibliography=totoc, listof=totoc]{scrreprt}

% --- Sprache & Encoding ---
\usepackage[utf8]{inputenc}
\usepackage[T1]{fontenc}
\usepackage[ngerman]{babel}
\usepackage{csquotes}

% --- Layout ---
\usepackage[left=3cm, right=2.5cm, top=2.5cm, bottom=2.5cm]{geometry}

% --- Literatur ---
\usepackage[backend=biber, style=authoryear]{biblatex}
\addbibresource{literatur.bib}

% --- Hyperlinks ---
\usepackage{hyperref}
\hypersetup{
    colorlinks=true,
    linkcolor=black,
    citecolor=black,
    urlcolor=blue
}

% --- Metadaten (Platzhalter) ---
\newcommand{\thesistitle}{Titel der Bachelorarbeit}
\newcommand{\authorname}{Vorname Nachname}
\newcommand{\matrikelnr}{Matrikelnummer}
\newcommand{\university}{Universität / Fachhochschule}
\newcommand{\faculty}{Fakultät / Studiengang}
\newcommand{\supervisor}{Betreuer/in: Prof.\ Dr.\ Vorname Nachname}
\newcommand{\secondreviewer}{Zweitgutachter/in: Prof.\ Dr.\ Vorname Nachname}
\newcommand{\submissiondate}{\today}

% ============================================================
\begin{document}

% --- Titelseite ---
\begin{titlepage}
    \centering
    \vspace*{2cm}
    {\Large \university \par}
    \vspace{0.5cm}
    {\large \faculty \par}
    \vspace{3cm}
    {\LARGE \textbf{\thesistitle} \par}
    \vspace{2cm}
    {\large Bachelorarbeit \par}
    \vspace{1cm}
    {\large vorgelegt von \par}
    \vspace{0.5cm}
    {\Large \authorname \par}
    {\normalsize Matrikelnr.: \matrikelnr \par}
    \vfill
    {\normalsize \supervisor \par}
    {\normalsize \secondreviewer \par}
    \vspace{1cm}
    {\normalsize \submissiondate \par}
\end{titlepage}

% --- Zusammenfassung ---
\chapter*{Zusammenfassung}
\addcontentsline{toc}{chapter}{Zusammenfassung}

% TODO: Zusammenfassung hier einfügen

% --- Inhaltsverzeichnis ---
\tableofcontents
\newpage

% ============================================================
% Hauptteil
% ============================================================

% Kapitel 0
\setcounter{chapter}{-1}
\chapter{Präambel: Zeitserie Pielach}

% Kapitel 1
\chapter{Einleitung}

% Kapitel 2
\chapter{Grundlagen und Stand der Technik}
\section{\textit{Cloud-Native-Geospatial} (COG, COPC und HTTP-Range-Requests)}
\section{Metadatenstandards: Von ISO~19115 zu STAC}
\section{Abgrenzung zum \textit{Data-Cube}-Paradigma}
\section{Verwandte Arbeiten und Werkzeuge}

% Kapitel 3
\chapter{Technische Ausgangslage der Rohdaten}
\section{Analyse der Verzeichnisstruktur und Dateiformate}
\section{Metadaten-Inkonsistenzen und programmatische Herausforderungen}

% Kapitel 4
\chapter{Konzeption der STAC-Architektur}
\section{Strukturierung der Katalog-Hierarchie}
\section{Evaluierung der \textit{Linking}-Strategie}
\section{Erweiterbarkeit und Idempotenz der Pipeline}
\section{Architektur der Abfrage- und \textit{Subsetting}-Komponente}

% Kapitel 5
\chapter{Implementierung}
\section{Datenaufbereitung und Überführung in \textit{Cloud-Native}-Formate}
\section{Generierungspipeline (Metadatenextraktion, konfigurierbare \textit{Asset-Registry})}
\section{Validierung (STAC-Konformität und inhaltliche Plausibilität)}
\section{Domänenspezifische Abfragen und \textit{Subsetting}}

% Kapitel 6
\chapter{Ergebnisse und Evaluierung}
\section{Präsentation des generierten STAC-Katalogs}
\section{Domänenspezifische raumzeitliche Abfragen und \textit{Subsetting}}
\section{Bewertung der Generalisierbarkeit der Pipeline}

% Kapitel 7
\chapter{Diskussion}
\section{Architekturentscheidungen: Begründung und Übertragbarkeit}
\section{Grenzen der Generalisierbarkeit und Konfigurierbarkeit}
\section{Mehrwert gegenüber verzeichnisorientierter Ablage und FAIR-Bezug}

% Kapitel 8
\chapter{Fazit und Ausblick}

% ============================================================
% Literaturverzeichnis
% ============================================================
\printbibliography

% ============================================================
% Anhang (optional)
% ============================================================
% \appendix
% \chapter{Anhang}

\end{document}
